\section{Conclusion}
\label{sec:conclusion}

The problem looked like a translation problem and was not. For Spanish the model-side
deficit is 1--4 points on reasoning and approximately zero on tool selection; the real
obstacle was that the system's own deterministic control plane was monolingual at the
tokenizer level, and that no amount of model quality could compensate for a user-chosen
model whose Spanish nobody had measured.

NEPANTLA resolves both by refusing to depend on the answer to the unanswerable question.
It \textbf{freezes} the user's literals before anything can touch them,
\textbf{neutralizes} the deterministic core so that intent --- not the language intent
arrived in --- drives routing, \textbf{proposes and verifies before it executes}, and
\textbf{escalates} through a ladder whose last rung is the English pipeline itself.
Because rejected proposals have no side effects, the ladder cannot end below its own
terminal rung; because the verifier's checks are decidable on exactly the failure classes
cross-lingual operation introduces, the false-accept rate on those classes is zero; and
because presentation is causally downstream of the verified action, the operator reads
Spanish without any of that reaching the machine.

Three consequences deserve to be carried away from this paper.

\textbf{Spanish competence became an optimization rather than a dependency.} A model with
no Spanish at all still produces correct execution and a Spanish answer; it merely costs
one or two extra calls. That is the direct answer to ``the user can pick any model''.

\textbf{The naming rule is a theorem, not a preference.} \code{Emailer} stays
\code{Emailer}, \code{Asker} stays \code{Asker}, \code{Apirer} stays \code{Apirer}, every
tool name, flag, configuration key, sentinel and identifier stays exactly as it is ---
because Proposition~1's independence, and therefore the entire guarantee, holds only
while the action vocabulary is the same set of symbols in every locale. The Spanish build
translates what the operator reads and nothing the machine reads. That line is drawn in
\S2.3 and enforced by Corollary~2.

\textbf{And the Spanish build may end up better than today's English one.} Three of the
required repairs --- boundary-aware phrase matching, the failure-classifier polarity, and
locale-invariant child processes --- fix defects that damage the English path right now.
Argument provenance protects an English user's paths exactly as well as a Spanish user's.
The honest stop at R4 prevents a fabricated deletion in either language.

Localization, done this way, is not a translation project attached to the side of a
product.

\begin{keyinsight}
It is a \textbf{correctness project that happened to be discovered by asking the system
to speak Spanish} --- and the discipline it forces, verify before you act and never let a
decision rest on a literal nobody said, is worth having regardless of the language anyone
is speaking.
\end{keyinsight}
